% Persona de coordinaci\'on de prevenci\'on
\begingroup
\fontsize{7.5}{9.7}\selectfont
\ifdefined\pcolw\else\newlength{\pcolw}\fi
\ifdefined\pfootw\else\newlength{\pfootw}\fi
\definecolor{PInk}{RGB}{55,48,163}
\definecolor{PTeal}{RGB}{13,148,136}
\definecolor{PBlue}{RGB}{37,99,235}
\definecolor{PSlate}{RGB}{71,85,105}
\definecolor{PEmr}{RGB}{5,150,105}
\definecolor{PAmb}{RGB}{217,119,6}
\definecolor{PRed}{RGB}{220,38,38}
\definecolor{PVio}{RGB}{124,58,237}
\definecolor{PInd}{RGB}{67,56,202}
\definecolor{PSky}{RGB}{2,132,199}
\definecolor{PPlum}{RGB}{190,24,93}
\newcommand{\PDot}[1]{%
  \tikz[baseline=-0.42ex]\fill[#1] (0,0) circle (1.85pt);%
  \hspace{0.48em}%
}
\newcommand{\PTit}[1]{\makebox[\linewidth][c]{\sffamily\bfseries #1}}
\tikzset{
  pchip/.style={
    sharp corners,
    rounded corners=7.7pt,
    inner xsep=6.8pt,
    inner ysep=0pt,
    minimum height=15.4pt,
    text height=0.68em,
    text depth=0.22em,
    font=\sffamily\bfseries\scriptsize,
    draw=none,
    line width=0pt,
    outer sep=0pt,
    align=center,
    anchor=base,
  },
}
\newcommand{\PPills}[6]{%
  \begin{tikzpicture}[baseline=0.22em, rounded corners=0pt, sharp corners]
    \node[pchip, fill=#1, text=#2] (p1) {#3};
    \node[pchip, fill=#4, text=#5, base right=4.2pt of p1] (p2) {#6};
  \end{tikzpicture}%
}
\newcommand{\PCaja}[4][south]{%
  \noindent
  \begin{tikzpicture}[baseline=(current bounding box.#1)]
    \begin{scope}[shift={(0.4pt,-1.8pt)}]
      \node[
        fill=black,
        opacity=0.18,
        text opacity=0,
        rounded corners=13pt,
        inner xsep=8.4pt,
        inner ysep=5.8pt,
        outer sep=0pt,
        align=left,
        font=\fontsize{7.5}{9.7}\selectfont\sffamily,
        text width=\dimexpr#2-16.8pt\relax,
      ] {#4};
    \end{scope}
    \node[
      fill=#3!9,
      rounded corners=13pt,
      inner xsep=8.4pt,
      inner ysep=5.8pt,
      outer sep=0pt,
      text=black,
      align=left,
      font=\fontsize{7.5}{9.7}\selectfont\sffamily,
      text width=\dimexpr#2-16.8pt\relax,
    ] (pcard) {#4};
    \fill[#3, rounded corners=2.1pt]
      ([xshift=2.2pt,yshift=5.2pt]pcard.south west)
      rectangle
      ([xshift=5.7pt,yshift=-5.2pt]pcard.north west);
  \end{tikzpicture}%
}

\begin{tikzpicture}
  \node[
    rounded corners=16pt,
    inner sep=0pt,
    outer sep=0pt,
    fill=black!3.5,
  ] {%
    \begin{minipage}{0.988\linewidth}
      \begin{tikzpicture}
        \node[
          fill=PInk,
          rounded corners=12pt,
          inner sep=0pt,
          outer sep=0pt,
          text width=\linewidth-1.4pt,
        ] {%
          \begin{minipage}{\linewidth}
            \vspace{7pt}
            \hspace{9pt}%
            \begin{minipage}{0.60\linewidth}
              \begin{tikzpicture}[baseline=-0.55ex]
                \fill[white] (0,0) circle (0.34cm);
                \node[font=\sffamily\bfseries\scriptsize, color=PInk] at (0,0) {LF};
              \end{tikzpicture}%
              \hspace{6pt}%
              \raisebox{0.15ex}{\textcolor{white}{\sffamily\bfseries\Large Laura Ferrero}}\\[3pt]
              \hspace{1pt}\textcolor{white}{\sffamily\scriptsize Un mapa de riesgo que diga d\'onde mirar, no un umbral que ya explot\'o}\\[2pt]
              \textcolor{white!88}{\sffamily\tiny Coordinaci\'on de prevenci\'on $\cdot$ ART $\cdot$ persona de gerencia}
            \end{minipage}%
            \hfill
            \begin{minipage}{0.34\linewidth}
              \raggedleft
              \PPills{white}{PInk}{Coordinadora}{PVio}{white}{Padr\'on y equipo}\\[5pt]
              \textcolor{white!88}{\sffamily\tiny Uso diario del tablero}
            \end{minipage}%
            \hspace{9pt}%
            \vspace{7pt}
          \end{minipage}%
        };
      \end{tikzpicture}

      \vspace{8pt}
      \raggedright
      \begin{minipage}{\dimexpr\linewidth-22pt\relax}
        \setlength{\pcolw}{\dimexpr(\linewidth-16pt)/3\relax}
        \setlength{\pfootw}{\dimexpr(\linewidth-8pt)/2\relax}

        \PCaja{\linewidth}{PTeal}{%
          {\color{PTeal}\PTit{Qui\'en es y c\'omo trabaja}}\\[4pt]
          Mira el padr\'on y decide d\'onde poner al equipo; no recorre cada visita.
          Eval\'ua la \enquote{salud} del establecimiento con el RGRL al afiliar o declarar, y el RAR/NTE para agentes de riesgo.
          El P.E.S.E.\ entra cuando el \'{\i}ndice de incidencia ya se dispar\'o: es un umbral reactivo.
          El plan de visitas vive en planillas; el sistema empuja m\'as a atender el siniestro que a prevenirlo.
          No denuncia todo: denunciar al afiliado deriva en inspectores y multa. El preventor verifica; la ART formaliza.
        }

        \vspace{9pt}
        \noindent
        \begin{minipage}[t]{\pcolw}
          \PCaja[north]{\pcolw}{PBlue}{%
            {\color{PBlue}\PTit{Contexto de uso}}\\[4pt]
            \PDot{PBlue}D\'onde: escritorio, tablero de cobertura y asignaci\'on.\\[2.6pt]
            \PDot{PBlue}Cu\'ando: arma el plan trimestral y rebalancea en la semana.\\[2.6pt]
            \PDot{PBlue}Por qu\'e: tiene que rendir capacidad de preventores, no un modelo.\\[2.6pt]
            \PDot{PBlue}Competencia: experta en RGRL/P.E.S.E., no en aprendizaje autom\'atico.
          }
        \end{minipage}%
        \hspace{8pt}%
        \begin{minipage}[t]{\pcolw}
          \PCaja[north]{\pcolw}{PSlate}{%
            {\color{PSlate}\PTit{Herramientas de hoy}}\\[4pt]
            \PDot{PSlate}RGRL como \enquote{biblia} de qu\'e prevenci\'on hacer.\\[2.6pt]
            \PDot{PSlate}Planillas para el plan de visitas.\\[2.6pt]
            \PDot{PSlate}\'Indice de incidencia publicado por la Superintendencia.\\[2.6pt]
            \PDot{PSlate}Sistema de gesti\'on de relevamientos y recomendaciones.
          }
        \end{minipage}%
        \hspace{8pt}%
        \begin{minipage}[t]{\pcolw}
          \PCaja[north]{\pcolw}{PEmr}{%
            {\color{PEmr}\PTit{Motivaciones}}\\[4pt]
            \PDot{PEmr}Ver el recorte accionable del padr\'on, no todo el universo.\\[2.6pt]
            \PDot{PEmr}Que el preventor sepa sector y agente material.\\[2.6pt]
            \PDot{PEmr}Cubrir incumplimientos con fecha de regularizaci\'on.\\[2.6pt]
            \PDot{PEmr}Equilibrar denuncia ante la SRT y relaci\'on con el afiliado.
          }
        \end{minipage}

        \vspace{9pt}
        \noindent
        \begin{minipage}[t]{\pcolw}
          \PCaja[north]{\pcolw}{PAmb}{%
            {\color{PAmb}\PTit{Objetivos}}\\[4pt]
            \PDot{PAmb}Distribuir preventores seg\'un riesgo relativo del trimestre.\\[2.6pt]
            \PDot{PAmb}Distinguir ranking predictivo de obligaci\'on PRS o P.E.S.E.\\[2.6pt]
            \PDot{PAmb}Medir cobertura y factores que m\'as empujan el ranking.\\[2.6pt]
            \PDot{PAmb}Sacar del foco del establecimiento los siniestros \emph{in itinere}.
          }
        \end{minipage}%
        \hspace{8pt}%
        \begin{minipage}[t]{\pcolw}
          \PCaja[north]{\pcolw}{PRed}{%
            {\color{PRed}\PTit{Frustraciones}}\\[4pt]
            \PDot{PRed}Actuar cuando el \'{\i}ndice ya estall\'o.\\[2.6pt]
            \PDot{PRed}Excel y RGRL manuales para priorizar gente.\\[2.6pt]
            \PDot{PRed}Poca inversi\'on relativa en prevenci\'on pese a los datos.\\[2.6pt]
            \PDot{PRed}\enquote{No se denuncia todo}: denunciar es denunciar al cliente.
          }
        \end{minipage}%
        \hspace{8pt}%
        \begin{minipage}[t]{\pcolw}
          \PCaja[north]{\pcolw}{PVio}{%
            {\color{PVio}\PTit{H\'abitos}}\\[4pt]
            \PDot{PVio}Pide al preventor el sector (poste, rampa) antes que un puntaje suelto.\\[2.6pt]
            \PDot{PVio}Cruza recomendaciones cumplidas o no con el RGRL.\\[2.6pt]
            \PDot{PVio}Separa accidentes de enfermedades profesionales (RAR).\\[2.6pt]
            \PDot{PVio}Quiere ver concentraci\'on anonimizada en el mismo trabajador.
          }
        \end{minipage}

        \vspace{10pt}
        \PCaja{\linewidth}{PInd}{%
          {\color{PInd}\PTit{Escenario}}\\[4pt]
          Abre estad\'{\i}sticas de cobertura del top-1000 y el SHAP agregado: d\'onde el modelo pone la masa.
          Asigna visitas a preventores sin mezclar un caso PRS con un cr\'{\i}tico predictivo.
          En la ficha mira incumplimientos y si la rampa del mes pasado ya se regulariz\'o; si est\'a cumplida, no la insiste.
          No trata el score como probabilidad. Afuera el trabajador casi no percibe prevenci\'on: el valor de Prevex es ordenar al equipo.
        }

        \vspace{9pt}
        \noindent
        \begin{minipage}[t]{\pfootw}
          \PCaja[north]{\pfootw}{PSky}{%
            {\color{PSky}\PTit{Qu\'e necesita de Prevex}}\\[4pt]
            \PDot{PSky}Cobertura y SHAP agregado sobre el top-1000 (HU-07, HU-09 / RF-07).\\[2.6pt]
            \PDot{PSky}Asignar visitas (HU-08 / RF-08).\\[2.6pt]
            \PDot{PSky}Distinguir motivo predictivo de PRS (HU-12).\\[2.6pt]
            \PDot{PSky}Ficha con historial e incumplimientos, sin leer el score como probabilidad (RF-04, RNF-06).
          }
        \end{minipage}%
        \hspace{8pt}%
        \begin{minipage}[t]{\pfootw}
          \PCaja[north]{\pfootw}{PPlum}{%
            {\color{PEmr}\PTit{\'Exito}}\\[3pt]
            El equipo mira primero donde hay que mirar, con causa accionable, y el plan no se mezcla con el umbral P.E.S.E.\\[5pt]
            {\color{PPlum}\PTit{Temor}}\\[3pt]
            Tratar el ranking como or\'aculo; ensuciar el establecimiento con \emph{in itinere}; romper la relaci\'on comercial con una denuncia autom\'atica.
          }
        \end{minipage}

        \vspace{4pt}
      \end{minipage}
    \end{minipage}%
  };
\end{tikzpicture}
\endgroup
